\textbf{Proof of (i):} We will show that $\expectationJ{\normtwo{\pi_{\partial J (\by)}(\boldsymbol{0})}} < +\infty$ and derive formulas~\eqref{eq:diff_mono_rep},~\eqref{eq:diff_pm_rep},~\eqref{eq:diff_var_rep},~\eqref{eq:diff_seps_rep}, and~\eqref{eq:diff_laplacian_rep} in four steps. To describe these steps, let us first introduce some notation. Recall that $J$ satisfies assumptions (A1)-(A3) and $\dom J = \Rn$. Define the set
\[
D_J\coloneqq\left\{\by \in \Rn \mid \partial J(\by) = \{\nabla J(\by)\} \right\}.
\] 
We can invoke (\cite{rockafellar1970convex}, Theorem 25.5)
to conclude that $D_J$ is a dense subset of $\Rn$, the $n$-dimensional Lebesgue measure of the set $(\Rn \setminus D_J)$ is zero, and the function $\by \mapsto \nabla J(\by)$ is continuous on $D_J$. Now, let $\bx \in \Rn$, $t>0$, $\epsilon > 0$, and $\by_0 \in \Rn$. Define the function $\varphi_J\colon \Rn \to \Rn$ as
\begin{equation*}
    \varphi_J(\by|\bx,t) = \left(\frac{\by-\bx}{t}\right) + \pi_{\partial J(\by)}(\boldsymbol{0}).
\end{equation*}
Note that for every $\by \in \Rn$ we have $\varphi_J(\by|\bx,t) \in \partial\left(
\Rn \ni \bu \mapsto \frac{1}{2t}\normsq{\bx-\bu} + J(\bu)
\right)(\by)$, i.e., $\varphi_J(\by|\bx,t)$ is a subgradient of the function $\bu \mapsto \frac{1}{2t}\normsq{\bx-\bu} + J(\bu)$ evaluated at $\bu = \by$. Let
\beq\label{eq:c1_appC}
C_1(\bx,\by_0,t,\epsilon) = \int_{\Rn} \normtwo{\by-\by_0}e^{-\frac{1}{2t\epsilon}\normsq{\bx-\by}}\diff{\by},
\eeq
and note that by assumption (A3), the expected value $\expectationJ{\normtwo{\by-\by_0}}$ is bounded as follows
\beq\label{eq:norm_appC}
\begin{alignedat}{1}
\expectationJ{\normtwo{\by-\by_0}} &= \frac{1}{Z_J(\bx,t,\epsilon)}\int_{\Rn}\normtwo{\by-\by_0}e^{-(\frac{1}{2t}\normsq{\bx-\by} + J(\by))/\epsilon}\diff{\by} \\
&\leqslant \frac{1}{Z_J(\bx,t,\epsilon)}\int_{\Rn}\normtwo{\by-\by_0}e^{-\frac{1}{2t\epsilon}\normsq{\bx-\by}}\diff{\by} \\
&= \frac{C_1(\bx,\by_0,t,\epsilon)}{Z_J(\bx,t,\epsilon)}.
\end{alignedat}
\eeq
Define the vector field $\bV\colon \Rn \to \Rn$ as
\beq \label{eq:V_field}
\bV(\by) = (\by-\by_0) e^{-\left(\frac{1}{2t}\left\Vert \bx-\by\right\Vert _{2}^{2} + J(\by)\right)/\epsilon},
\eeq
which is continuous on $\Rn$. It is also bounded on $\Rn$; to see this, use the triangle inequality, assumption (A3), and the fact that the function $(0,+\infty) \ni r \mapsto re^{-\frac{1}{2t\epsilon}r^2}$ attains its maximum at $r^* = \sqrt{t\epsilon}$ to get
\beq\label{eq:boundV1_appC}
\begin{alignedat}{1}
\normtwo{\bV(\by)} &= \normtwo{(\by-\by_0)} e^{-\left(\frac{1}{2t}\left\Vert \bx-\by\right\Vert _{2}^{2} + J(\by)\right)/\epsilon} \\
&=\normtwo{(\by-\bx + \bx-\by_0)} e^{-\left(\frac{1}{2t}\left\Vert \bx-\by\right\Vert _{2}^{2} + J(\by)\right)/\epsilon} \\
&\leqslant (\normtwo{\bx-\by_0} + \normtwo{\bx-\by}) e^{-\left(\frac{1}{2t}\left\Vert \bx-\by\right\Vert _{2}^{2} + J(\by)\right)/\epsilon} \\
&\leqslant \normtwo{\bx-\by_0} + \normtwo{\bx-\by}e^{-\left(\frac{1}{2t}\left\Vert \bx-\by\right\Vert _{2}^{2} + J(\by)\right)/\epsilon} \\
&\leqslant \normtwo{\bx-\by_0} + \normtwo{\bx-\by}e^{-\frac{1}{2t\epsilon}\left\Vert \bx-\by\right\Vert _{2}^{2}} \\
&\leqslant  \normtwo{\bx-\by_0} + \sup_{\by\in\Rn}\left( \normtwo{\bx-\by}e^{-\frac{1}{2t\epsilon}\left\Vert \bx-\by\right\Vert _{2}^{2}}\right) \\
&\leqslant \normtwo{\bx-\by_0} +  (\sqrt{t\epsilon})e^{-\frac{\sqrt{t\epsilon}}{2}}.
\end{alignedat}
\eeq
The divergence $\nabla_{\by} \cdot \bV(\by)$, which is well-defined and continuous on $D_J$, is given for every $\by \in D_J$ by
\beq \label{eq:div_V}
\begin{alignedat}{1}
\nabla_{\by}\cdot\bV(\by) &= \nabla_{\by}\cdot\left((\by-\by_0) e^{-\left(\frac{1}{2t}\normsq{\bx-\by} + J(\by)\right)/\epsilon}\right) \\
&= (\nabla_{\by}\cdot (\by-\by_0))e^{-\left(\frac{1}{2t}\left\Vert \bx-\by\right\Vert _{2}^{2} + J(\by)\right)/\epsilon} + \left\langle \nabla_{\by}e^{-\left(\frac{1}{2t}\left\Vert \bx-\by\right\Vert _{2}^{2} + J(\by)\right)/\epsilon}, \by-\by_0 \right\rangle \\
&= ne^{-\left(\frac{1}{2t}\left\Vert \bx-\by\right\Vert _{2}^{2} + J(\by)\right)/\epsilon} - \left\langle \frac{1}{\epsilon}\left(\frac{\by - \bx}{t} + \nabla J(\by)\right)e^{-\left(\frac{1}{2t}\left\Vert \bx-\by\right\Vert _{2}^{2} + J(\by)\right)/\epsilon}, \by-\by_0 \right\rangle\\
&=\left(n  - \left\langle \frac{1}{\epsilon}\left(\frac{\by - \bx}{t} + \nabla J(\by)\right), \by-\by_0 \right\rangle \right)e^{-\left(\frac{1}{2t}\left\Vert \bx-\by\right\Vert _{2}^{2} + J(\by)\right)/\epsilon}.
\end{alignedat}
\eeq

We now outline the four steps that will be used to prove Prop.~\ref{prop:pm_rep_props}(i). In the first step, we will show that the divergence of the vector field $\bV$ on $D_J$ integrates to zero in the sense that
\beq\label{eq:appC_int}
\lim_{r \to +\infty} \left|\int_{\left\{\by \in \Rn \mid\, \normtwo{\by} \leqslant r\right\}\cap\,D_J} \nabla_{\by} \cdot \bV(\by) \diff{\by}\right| = 0.
\eeq
In the second step, we will show that $\expectationJ{\left|\left\langle \varphi_J(\by|\bx,t),\by-\by_0 \right\rangle\right|} < +\infty$, with $\expectationJ{\left\langle \varphi_J(\by|\bx,t),\by-\by_0 \right\rangle} = n\epsilon$, hereby proving formula~\eqref{eq:diff_mono_rep}, using the convexity of the function $\by \mapsto \frac{1}{2t}\normsq{\bx-\by} + J(\by)$, Fatou's lemma (\citep{folland2013real}, Lemma 2.18), and Eq.~\eqref{eq:appC_int} derived in the first step. In the third step, we will combine the results from the first and second steps to show that $\expectationJ{\normtwo{\pi_{\partial J (\by)}(\boldsymbol{0})}} < +\infty$ and conclude that the representation formulas~\eqref{eq:diff_pm_rep} and~\eqref{eq:diff_var_rep} hold. Finally, in the fourth step we will conclude that the representation formulas~\eqref{eq:diff_seps_rep} and~\eqref{eq:diff_laplacian_rep} hold using Eqs.~\eqref{eq:diff_pm_rep} and~\eqref{eq:diff_var_rep} and Prop.~\eqref{thm:visc_hj_eqns}(iii).


%%%%%%%%%%%%%%%%%%%%%%%%%%%%%%%%
\bigbreak
\noindent
\textbf{Step 1.} The proof of the limit result~\eqref{eq:appC_int} that we present here is based on an application of Theorem 4.14 in \cite{pfeffer1990divergence} to the vector field $\bV(\cdot)$. As this result is fairly technical, we first introduce some terminology and definitions that will be used exclusively in this part of the proof of (i).

\bigbreak
\noindent
Let $C$ be a non-empty convex subset of $\Rn$. The \textit{dimension} of the set $C$ is defined as the smallest dimension of a non-empty affine set containing $C$, with the dimension of a non-empty affine set being the dimension of the subspace parallel to it (\cite{rockafellar1970convex}, pages 4 and 12). If $C$ consists of a single point then its dimension is taken to be zero.

\bigbreak
\noindent
Let $k \in \{0,\dots,n\}$. Denote by $\mathcal{H}^{n-k}$ the $(n-k)$-dimensional outer Hausdorff measure in $\Rn$ as defined in (\cite{federer1969geometric}, Sect. 2.10.2, p.171). The measure $\mathcal{H}^{n-k}$, in particular, is a constant multiple of the $(n-k)$-dimensional Lebesgue measure for every measurable subset $B \subset \Rn$ (see \cite{falconer_1985}, Section 1.2, p.7, and Theorem 1.12, p.13). 

\bigbreak
\noindent
A subset $S \subset \Rn$ is called \textit{slight} if $\mathcal{H}^{n-1}(S) = 0$, and a subset $T \subset \Rn$ is called \textit{thin} if $T$ is $\sigma$-finite for $\mathcal{H}^{n-1}$, i.e., $T$ can be expressed as a countable union of sets $T = \cup_{k=1}^{+\infty} T_k$ with $\mathcal{H}^{n-1}(T_k) < +\infty$ for each $k \in \mathbb{N}^+$ (see, e.g., \cite{pfeffer1990divergence}).

\bigbreak
\noindent
Let $k \in \{0,\dots,n\}$. A non-empty, measurable subset $\Omega \subset \Rn$ is said to be \textit{countably $\mathcal{H}^{n-k}$-rectifiable} if it is contained, up to a null set of $(n-k)$-dimensional outer Hausdorff measure $\mathcal{H}^{n-k}$ zero, in a countable union of continuously differentiable hypersurfaces of dimension $(n-k)$ (see, e.g., \cite{alberti1992singularities} and references therein). A non-empty, measurable and countably $\mathcal{H}^{n-k}$-rectifiable subset of $\Rn$, in particular, is $\sigma$-finite for $\mathcal{H}^{n-k}$. 

\bigbreak
\noindent
A subset $A\subset \Rn$ is called \textit{admissible} if its boundary $\bd A$ is thin and if the distributional gradient of the characteristic function of $A$ is a vector measure on Borel subsets of $\Rn$ whose variation is finite (see \cite{pfeffer1990divergence} pp.151 and the reference therein). For the purpose of our proof, we will use the fact that the family of closed balls of radius $r>0$, namely $\{\by \in \Rn \mid \normtwo{\by} \leqslant r\}$, are admissible sets (see \cite{giusti1984minimal}, Example 1.10, and note that admissible sets are also called \textit{Caccioppoli} sets \cite{giusti1984minimal}, pages 5-6).

\bigbreak
\noindent
Let $A$ be an admissible set and let $\bv\colon A\to\Rn$ be a vector field. In the terminology of \cite{pfeffer1990divergence}, we say that $\bv$ is \textit{integrable} over the admissible set $A$ if $\bv$ satisfies definition 4.1 of \cite{pfeffer1990divergence}, and in that case, the number $I(\bv,A)$ is called the integral of $\bv$ over $A$. Note, here, that the notion of integrability considered in $\cite{pfeffer1990divergence}$ is different from that of the usual Lebesgue integrability. Nevertheless, if $\bv$ is integrable in the sense of \cite{pfeffer1990divergence}, $\bv$ is also Lebesgue measurable (Corollary 4.9, \cite{pfeffer1990divergence}), and if the Lebesgue integral $\int_{A}|\bv(\by)|\diff{\by}$ is finite, then $I(\bv,A) = \int_{A}|\bv(\by)|\diff{\by}$ (\cite{pfeffer1990divergence}, Prop. 4.7).
 
\bigbreak
\noindent
Let $E$ be an non-empty subset of $\Rn$, let $\bv\colon E \to \Rn$ be a vector field, and let $D_{\bv}$ denote the set of points at which $\bv$ is differentiable in $\inter{E}$ (for the definition of differentiability of vector fields, see \cite{rudin2006real}, page 150, Definition 7.22). In the terminology of \cite{pfeffer1990divergence}, we call a \textit{divergence of $\bv$} any function $g \colon E \mapsto \R$ such that $g(\by) = \nabla \cdot \bv(\by)$ for each $\by \in (\inter{E}) \cap D_{\bv}$.

\bigbreak
\noindent
In addition to these definitions, we will need the following two results due to, respectively, \cite{alberti1992singularities} and \cite{pfeffer1990divergence}.
\begin{thm}\label{thm:alberti}
[\cite{alberti1992singularities}, Theorem 4.1 (for convex functions)]
Let $\Omega$ be a bounded, open, convex subset of $\Rn$, and let $f\colon\Omega\to\R$ be a convex and Lipschitz continuous function. Denote the subdifferential of $f$ at $\by \in \Omega$ by $\partial f(\by)$. Then, for each $k\in\{0,\dots,n\}$, the set
\[
\{\by \in \Omega \mid \text{dim}\left(\partial f(\by)\right)\geqslant k\}
\]
is countably $\mathcal{H}^{n-k}$-rectifiable.
\end{thm}

\begin{thm}\label{thm:pfeffer}
[\cite{pfeffer1990divergence}, Theorem 4.14]
Let $A$ be an admissible set, and let $S$ and $T$ be, respectively, a slight and thin subset of $\cl{A}$. Let $\bv$ be a bounded vector field in $\cl{A}$ that is continuous in $(\cl{A})\setminus S$ and differentiable in $(\inter{A})\setminus T$. Then every divergence of $\bv$ is integrable in $A$. Moreover, there exists a vector field $\bd A \ni \by \to n_{A}(\by)$ with $\normtwo{n_{A}(\by)} = 1$ for every $\by \in \bd A$ such that if $\text{div }\bv$ denotes any divergence $\bv$, then
\beq
I(\text{div }\bv,A) = \int_{\bd{A}} \left\langle\bv(\by), \bn_{\bv}(\by)\right\rangle d\mathcal{H}^{n-1}\diff{\by}.
\eeq
\end{thm}

\bigbreak
\noindent
\textbf{Step 1 (Continued).} Fix $r>0$ and let $A = \{\by \in \Rn \mid \normtwo{\by} \leqslant r\}$ denote the closed ball of radius $r$ centered at the origin in $\Rn$. Note that $A$ is bounded, convex, closed, and admissible. Consider now the restriction of the convex function $J$ to $\inter{A}$. As $\inter{A}$ is bounded, open and convex, the function $J$ is Lipschitz continuous on $\inter{A}$ (\cite{rockafellar1970convex}, Theorem 10.4). All conditions in Thm.~\eqref{thm:alberti} are satisfied (with $\Omega = \inter{A}$ and $f = J$), and we can invoke the theorem to conclude that the set 
\[
T = \{\by \in \inter{A} \mid \text{dim}\left(\partial J(\by)\right)\geqslant 1\}
\]
is countably $\mathcal{H}^{n-1}$-rectifiable, and therefore $\sigma$-finite for $\mathcal{H}^{n-1}$. In particular, the set $T$ is thin. Moreover, recalling the definition of the set $D_J\coloneqq\left\{\by \in \Rn \mid \partial J(\by) = \{\nabla J(\by)\} \right\}$, we find that the set $(\inter{A})\setminus T$ comprises the points $\by \in \inter{A}$ at which the subdifferential $\partial J(\by)$ is a singleton, i.e., $T = (\inter{A})\cap (\Rn\setminus D_J)$.

Now consider the vector field $\bV$ defined by~\eqref{eq:V_field}. This vector field is continuous in $\Rn$ by convexity of $J$ and $\dom J = \Rn$. It is also bounded by~\eqref{eq:boundV1_appC}. Now define the function $g \colon A \to \R$ via
\beq\label{eq:appC_div1}
g(\by) = 
\begin{cases} 
\nabla \cdot \bV(\by) \text{ if }\by\in A\cap D_{J}, \\ 
0, \quad\text{ if } \by\in A\cap(\Rn\setminus D_{J}).
\end{cases}
\eeq
The function $g$ constitutes a divergence of the vector field $\bV$ because it coincides with the divergence $\nabla \cdot \bV(\by)$ at every $\by \in (\inter{A})\cap D_J$. Moreover, its Lebesgue integral over $A$ is finite; to see this, first note that for every $\by \in A\cap D_J$ the absolute value of $g(\by)$ can be bounded using~\eqref{eq:div_V}, the triangle inequality, the Cauchy--Schwarz inequality, and assumption (A3) as follows
\beq\label{eq:div2_appD}
\begin{alignedat}{1}
|g(\by)| = \left| \nabla \cdot \bV(\by) \right| &= \left|\left(n  - \left\langle \frac{1}{\epsilon}\left(\frac{\by - \bx}{t} + \nabla J(\by)\right), \by-\by_0 \right\rangle \right)e^{-\left(\frac{1}{2t}\left\Vert \bx-\by\right\Vert _{2}^{2} + J(\by)\right)/\epsilon}\right| \\
&\leqslant \left(n + \left|\left\langle \frac{1}{\epsilon}\left(\frac{\by - \bx}{t} + \nabla J(\by)\right), \by-\by_0 \right\rangle \right|\right)e^{-\left(\frac{1}{2t}\left\Vert \bx-\by\right\Vert _{2}^{2} + J(\by)\right)/\epsilon} \\
&\leqslant \left(n +  \frac{1}{\epsilon}\normtwo{\frac{\by - \bx}{t} + \nabla J(\by)}\normtwo{\by-\by_0} \right)e^{-\left(\frac{1}{2t}\left\Vert \bx-\by\right\Vert _{2}^{2} + J(\by)\right)/\epsilon} \\
&\leqslant \left(n +  \frac{1}{\epsilon}\left(\normtwo{\frac{\by - \bx}{t}} + \normtwo{\nabla J(\by)}\right)\normtwo{\by-\by_0} \right)e^{-\left(\frac{1}{2t}\left\Vert \bx-\by\right\Vert _{2}^{2} + J(\by)\right)/\epsilon} \\ 
&\leqslant \left(n +  \frac{1}{\epsilon}\left(\normtwo{\frac{\by - \bx}{t}} + \normtwo{\nabla J(\by)}\right)\normtwo{\by-\by_0} \right)e^{-\frac{1}{2t\epsilon}\left\Vert \bx-\by\right\Vert _{2}^{2}}
\end{alignedat}
\eeq
Second, as the set A is a closed bounded subset of $\dom J = \Rn$ the function $J$ is Lipschitz continuous relative to $A$, and therefore there exists a number $L_A>0$ such that $\normtwo{\nabla J(\by)} \leqslant L_A$ for every $\by \in A \cap D_J$. As a consequence, we can further bound $g(\by)$ for every $\by \in A\cap D_J$ in~\eqref{eq:div2_appD} as
\beq\label{eq:div3_appD}
|g(\by)| \leqslant \left(n +  \frac{1}{\epsilon}\left(\normtwo{\frac{\by - \bx}{t}} + L_A\right)\normtwo{\by-\by_0} \right)e^{-\frac{1}{2t\epsilon}\left\Vert \bx-\by\right\Vert _{2}^{2}}.
\eeq
In particular, using the definition of $g$ given by~\eqref{eq:appC_div1}, we have that~\eqref{eq:div3_appD} holds for every $\by \in A$. We can now use~\eqref{eq:appC_div1} and~\eqref{eq:div3_appD} to get
\beq\label{eq:div4_appD}
\begin{alignedat}{1}
\int_{A}|g(\by)|\diff \by = \int_{\{\by \in \Rn \mid \normtwo{\by} \leqslant r\}}|g(\by)|\diff\by &= \int_{\{\by \in \Rn \mid \normtwo{\by} \leqslant r\}\cap D_J} \left| \nabla \cdot \bV(\by) \right|\diff\by \\
&\leqslant \int_{\{\by \in \Rn \mid \normtwo{\by} \leqslant r\}\cap D_J} \left(n +  \frac{1}{\epsilon}\left(\normtwo{\frac{\by - \bx}{t}} + L_A\right)\normtwo{\by-\by_0} \right)e^{-\frac{1}{2t\epsilon}\left\Vert \bx-\by\right\Vert _{2}^{2}} \diff\by \\
&\leqslant \int_{\{\by \in \Rn \mid \normtwo{\by} \leqslant r\}} \left(n +  \frac{1}{\epsilon}\left(\normtwo{\frac{\by - \bx}{t}} + L_A\right)\normtwo{\by-\by_0} \right)e^{-\frac{1}{2t\epsilon}\left\Vert \bx-\by\right\Vert _{2}^{2}} \diff\by. \\
\end{alignedat}
\eeq
Since the function
\[
\by \mapsto \left(n +  \frac{1}{\epsilon}\left(\normtwo{\frac{\by - \bx}{t}} + L_A\right)\normtwo{\by-\by_0} \right)e^{-\frac{1}{2t\epsilon}\left\Vert \bx-\by\right\Vert _{2}^{2}}
\]
is continuous, it is bounded on the compact set $A = \{\by \in \Rn \mid \normtwo{\by} \leqslant r\}$. Its integral over $A$ is therefore finite, and using~\eqref{eq:div4_appD} we find that $\int_{A}|g(\by)|\diff \by$ is finite as well.

The previous considerations show that all conditions in Thm.~\eqref{thm:pfeffer} are satisfied (with $A = \{\by \in \Rn \mid \normtwo{\by} \leqslant r\}$, $\bv = \bV$, $S=\varnothing$, $T = \{\by \in \inter{A} \mid \text{dim}\left(\partial J(\by)\right)\geqslant 1\} = (\inter{A})\cap (\Rn\setminus D_J)$). We can therefore invoke (\cite{pfeffer1990divergence}, Theorem 4.14) to conclude that the divergence of $g$ is integrable (in the sense described by \cite{pfeffer1990divergence}), with integral $I(g,A)$, and that there exists a vector field $\bd A \ni \by \to \bn_{\bv}(\by)$ with $\normtwo{\bn_{\bv}(\by)} = 1$ for every $\by \in \bd A$ such that 
\beq\label{eq:appC_div2}
I(g,A) = \int_{\bd A} \left\langle\bV(\by), \bn_{\bv}(\by)\right\rangle \, d\mathcal{H}^{n-1}\diff{\by}.
\eeq
Since the Lebesgue integral of $|g|$ over $A$ is finite, we also have (\cite{pfeffer1990divergence}, Prop 4.7)
\beq\label{eq:iga_appC}
I(g,A) = \int_{A} g(\by) \diff\by.
\eeq
Using that $A = \{\by \in \Rn \mid \normtwo{\by} \leqslant r\}$, Eqs.~\eqref{eq:appC_div1},~\eqref{eq:appC_div2}, and~\eqref{eq:iga_appC}, we obtain
\beq\label{eq:appC_div4}
\int_{\{\by \in \Rn \mid \normtwo{\by} \leqslant r\}} g(\by) \diff\by = \int_{\{\by \in \Rn \mid \normtwo{\by} \leqslant r\}} \nabla\cdot\bV(\by) \diff\by = \int_{\{\by \in \Rn \mid \normtwo{\by} = r\}} \left\langle\bV(\by), \bn_{\bv}(\by)\right\rangle \, d\mathcal{H}^{n-1}.
\eeq
As $r$ was an arbitrary positive number, we can take the absolute value and then the limit $r \to +\infty$ on both sides of~\eqref{eq:appC_div4} to find
\beq\label{eq:appC_div5}
\lim_{r\to+\infty} \left|\int_{\{\by \in \Rn \mid \normtwo{\by} \leqslant r\}\cap D_J} \nabla\cdot\bV(\by) \diff\by \right| = \lim_{r\to +\infty} \left| \int_{\{\by \in \Rn \mid \normtwo{\by} = r\}} \left\langle\bV(\by), \bn_{\bv}(\by)\right\rangle \, d\mathcal{H}^{n-1} \right|.
\eeq

We will now show that the limit on the right side of~\eqref{eq:appC_div5} is equal to zero. To show this, first take the absolute value inside the integral on the right side of~\eqref{eq:appC_div5} to find
\beq\label{eq:appC_div6}
\begin{alignedat}{1}
\left| \int_{\{\by \in \Rn \mid \normtwo{\by} = r\}} \left\langle \bV(\by), \bn_{\bv}(\by)\right\rangle \, d\mathcal{H}^{n-1} \right| &\leqslant \int_{\{\by \in \Rn \mid \normtwo{\by} = r\}} \left|\left\langle\bV(\by), \bn_{\bv}(\by)\right\rangle\right| \, d\mathcal{H}^{n-1}.
\end{alignedat}
\eeq
Use the Cauchy--Schwarz inequality, Eq.~\eqref{eq:V_field}, assumption (A3) ($\inf_{\by\in\Rn}J(\by)=0$) and $\normtwo{\bn_{\bv}} = 1$ to further bound the right side of~\eqref{eq:appC_div6} as follows
\beq\label{eq:appC_div7}
\begin{alignedat}{1}
\int_{\{\by \in \Rn \mid \normtwo{\by} = r\}} \left|\left\langle\bV(\by), \bn_{\bv}(\by)\right\rangle\right| \, d\mathcal{H}^{n-1} &\leqslant \int_{\{\by \in \Rn \mid \normtwo{\by} = r\}} \normtwo{\bV(\by)}\normtwo{\bn_{\bv}(\by)} \, d\mathcal{H}^{n-1} \\
&\leqslant \int_{\{\by \in \Rn \mid \normtwo{\by} = r\}} \normtwo{\by-\by_0} e^{-\left(\frac{1}{2t}\left\Vert \bx-\by\right\Vert _{2}^{2} + J(\by)\right)/\epsilon} \, d\mathcal{H}^{n-1}\\
&\leqslant \int_{\{\by \in \Rn \mid \normtwo{\by} = r\}} (\normtwo{\by} + \normtwo{\by_0}) e^{-\left(\frac{1}{2t}\left\Vert \bx-\by\right\Vert _{2}^{2}\right)/\epsilon} \, d\mathcal{H}^{n-1}.
\end{alignedat}
\eeq
Use the parallelogram law $2(\normsq{\bx} + \normsq{\bu}) = \normsq{\bx-\bu} + \normsq{\bx+\bu}$ with $\bu = \bx-\by$ to bound the exponential $e^{-\frac{1}{2t}\normsq{\bx - \by}/\epsilon}$ by
\beq\label{eq:appC_exp_term}
\begin{alignedat}{1}
e^{-\frac{1}{2t}\normsq{\bx - \by}/\epsilon} &= e^{-\frac{1}{2t}(\frac{1}{2}(\normsq{\by} + \normsq{2\bx-\by}) - \normsq{\bx})/\epsilon} \\
&\leqslant e^{-\frac{1}{2t}(\frac{1}{2}\normsq{\by} - \normsq{\bx})/\epsilon}
\end{alignedat}
\eeq
and use it in~\eqref{eq:appC_div7} to get
\beq\label{eq:appC_div8}
\int_{\{\by \in \Rn \mid \normtwo{\by} = r\}} \left|\left\langle\bV(\by), \bn_{\bv}(\by)\right\rangle\right| \, d\mathcal{H}^{n-1} \leqslant \int_{\{\by \in \Rn \mid \normtwo{\by} = r\}} (\normtwo{\by} + \normtwo{\by_0}) e^{-\frac{1}{2t}(\frac{1}{2}\normsq{\by} - \normsq{\bx})/\epsilon}  \, d\mathcal{H}^{n-1}.
\eeq
Since the domain of integration in~\eqref{eq:appC_div8} is over the surface of an $n$-dimensional sphere of radius $\normtwo{\by} = r$, the integral on the right side of~\eqref{eq:appC_div8} is given by
\beq\label{eq:appC_div9}
\begin{alignedat}{1}
\int_{\{\by \in \Rn \mid \normtwo{\by} = r\}} (\normtwo{\by} + \normtwo{\by_0}) e^{-\frac{1}{2t}(\frac{1}{2}\normsq{\by} - \normsq{\bx})/\epsilon}  \, d\mathcal{H}^{n-1} &= \int_{\{\by \in \Rn \mid \normtwo{\by} = r\}} (r + \normtwo{\by_0}) e^{-\frac{1}{2t}(\frac{1}{2}r^2 - \normsq{\bx})/\epsilon}  \, d\mathcal{H}^{n-1} \\
&= (r + \normtwo{\by_0}) e^{-\frac{1}{2t}(\frac{1}{2}r^2 - \normsq{\bx})/\epsilon} \int_{\{\by \in \Rn \mid \normtwo{\by} = r\}} d\mathcal{H}^{n-1} \\
&=(r + \normtwo{\by_0})e^{-\frac{1}{2t}\left(\frac{1}{2}r^2 - \normsq{\bx} \right)/\epsilon} \frac{n\pi^{n/2}}{\Gamma\left(\frac{n}{2} + 1\right)},
\end{alignedat}
\eeq
where $n\pi^{n/2}/\Gamma\left(\frac{n}{2} + 1\right)$ is the area of an $n$-dimensional sphere of radius one, with
$\Gamma\left(\frac{n}{2} + 1\right)$ denoting the Gamma function evaluated at $\frac{n}{2} + 1$. Since 
\[
\lim_{r\to+\infty} (r + \normtwo{\by_0})e^{-\frac{1}{2t}\left(\frac{1}{2}r^2 - \normsq{\bx} \right)/\epsilon} = 0,
\]
the limit $r\to+\infty$ in~\eqref{eq:appC_div9} is equal to zero, i.e.,
\beq\label{eq:appC_div10}
\lim_{r\to +\infty}\int_{\{\by \in \Rn \mid \normtwo{\by} = r\}} (\normtwo{\by} + \normtwo{\by_0}) e^{-\frac{1}{2t}(\frac{1}{2}\normsq{\by} - \normsq{\bx})/\epsilon}  \, d\mathcal{H}^{n-1} = 0.
\eeq
Combining~\eqref{eq:appC_div5},~\eqref{eq:appC_div6},~\eqref{eq:appC_div8} and~\eqref{eq:appC_div10} yield
\[
\lim_{r \to +\infty} \left|\int_{\left\{\by \in \Rn \mid\, \normtwo{\by} \leqslant r\right\}\cap\,D_J} \nabla_{\by} \cdot \bV(\by) \diff{\by}\right| = 0.
\]
which proves the limit result~\eqref{eq:appC_int}.

%%%%%%%%%%%%%%%%%%%%%%

\bigbreak
\noindent
\textbf{Step 2.} 
Recall that the divergence of the vector field $\by \mapsto \bV(\by)$ on $D_J$ is given by~\eqref{eq:div_V}. Combine~\eqref{eq:appC_int} and~\eqref{eq:div_V} to conclude that
\beq \label{eq:first_int_appC}
\lim_{r \to +\infty} \left|\int_{\left\{\by \in \Rn \mid\, \normtwo{\by} \leqslant r\right\}\cap\,D_J} \left(n\epsilon  - \left\langle \left(\frac{\by - \bx}{t} + \nabla J(\by)\right), \by-\by_0 \right\rangle \right)e^{-\left(\frac{1}{2t}\left\Vert \bx-\by\right\Vert _{2}^{2} + J(\by)\right)/\epsilon} \diff{\by}\right| = 0.
\eeq
Note that the minimal subgradient $\pi_{\partial J(\by)}(\boldsymbol{0}) = \nabla J(\by)$ for every $\by \in D_J$. We can therefore substitute the minimal subgradient $\pi_{\partial J(\by)}(\boldsymbol{0})$ for the gradient $\nabla J(\by)$ inside the integral in the limit~\eqref{eq:first_int_appC} without changing its value. Moreover, since the set $D_J$ is dense in $\Rn$ and the $n$-dimensional Lebesgue measure of $(\Rn \setminus D_J)$ is zero, we can further substitute the domain of integration $\left\{\by \in \Rn \mid\, \normtwo{\by} \leqslant r\right\}\cap\,D_J$ of the integral in the limit~\eqref{eq:first_int_appC} with $\left\{\by \in \Rn \mid\, \normtwo{\by} \leqslant r\right\}$ without changing its value. With these two changes, the limit~\eqref{eq:first_int_appC} can be written as
\[
\lim_{r \to +\infty} \left|\int_{\left\{\by \in \Rn \mid\, \normtwo{\by} \leqslant r\right\}} \left(n\epsilon  - \left\langle \left(\frac{\by - \bx}{t} + \pi_{\partial J(\by)}(\boldsymbol{0})\right), \by-\by_0 \right\rangle \right)e^{-\left(\frac{1}{2t}\left\Vert \bx-\by\right\Vert _{2}^{2} + J(\by)\right)/\epsilon} \diff{\by}\right| = 0.
\]
Using the notation $\varphi_J(\by|\bx,t) = \left(\frac{\by - \bx}{t}\right) + \pi_{\partial J(\by)}(\boldsymbol{0})$, we can write this limit more succinctly as
\beq\label{eq:first_int2_appC}
\lim_{r \to +\infty} \left|\int_{\left\{\by \in \Rn \mid\, \normtwo{\by} \leqslant r\right\}} \left(n\epsilon  - \left\langle \varphi_J(\by|\bx,t), \by-\by_0 \right\rangle \right)e^{-\left(\frac{1}{2t}\left\Vert \bx-\by\right\Vert _{2}^{2} + J(\by)\right)/\epsilon} \diff{\by}\right| = 0.
\eeq
Now, consider the function $\Rn \ni \by \mapsto \left\langle \varphi_J(\by|\bx,t) - \varphi_J(\by_0|\bx,t),\by-\by_0 \right\rangle$. Note here that as $J$ is convex with $\dom J =\Rn$, both $\varphi_J(\by|\bx,t)$ and $\varphi_J(\by_0|\bx,t)$ are subgradients of the convex function $\bu \mapsto \frac{1}{2t}\normsq{\bx-\bu} + J(\bu)$ at $\bu = \by$ and $\bu = \by_0$, respectively (\citep{rockafellar1970convex}, Theorem 23.4). We can therefore apply inequality~\eqref{eq:monotone_mapping} (with $p = \varphi_J(\by|\bx,t)$, $p_0 = \varphi_J(\by_0|\bx,t)$, and $m = 0$) to find $\left\langle \varphi_J(\by|\bx,t) - \varphi_J(\by_0|\bx,t),\by-\by_0 \right\rangle \geqslant 0$. Define $F\colon \Rn \to \R$ and $G \colon \Rn \to \R$ as follows:
\[
F(\by) = \left\langle \varphi_J(\by|\bx,t),\by-\by_0 \right\rangle e^{-\left(\frac{1}{2t}\left\Vert \bx-\by\right\Vert _{2}^{2} + J(\by)\right)/\epsilon}
\]
and
\[
G(\by) = \left\langle \varphi_J(\by_0|\bx,t),\by-\by_0 \right\rangle e^{-\left(\frac{1}{2t}\left\Vert \bx-\by\right\Vert _{2}^{2} + J(\by)\right)/\epsilon}.
\]
Note that $F(\by)-G(\by) = \left\langle \varphi_J(\by|\bx,t) - \varphi_J(\by_0|\bx,t),\by-\by_0 \right\rangle e^{-\left(\frac{1}{2t}\left\Vert \bx-\by\right\Vert _{2}^{2} + J(\by)\right)/\epsilon} \geqslant 0$ for every $\by \in \Rn$. Integrate $\by \mapsto F(\by)-G(\by)$ over $\Rn$ and use Fatou's lemma to find
\beq\label{eq:first_int3_appC}
\begin{alignedat}{1}
0 \leqslant \int_{\Rn} F(\by)-G(\by) \diff{\by} &\leqslant \lim_{r \to +\infty} \int_{\left\{\by \in \Rn \mid\, \normtwo{\by} \leqslant r\right\}} F(\by)-G(\by)  \diff{\by} \\
&= \lim_{r \to +\infty} \left(\int_{\left\{\by \in \Rn \mid\, \normtwo{\by} \leqslant r\right\}} F(\by) \diff{\by} + \int_{\left\{\by \in \Rn \mid\, \normtwo{\by} \leqslant r\right\}} (-G(\by))  \diff{\by}\right)
\end{alignedat}
\eeq
Use the Cauchy--Schwarz inequality assumption (A3) ($\inf_{\by\in\Rn} J(\by) = 0$) to bound the second integral on the right hand side of~\eqref{eq:first_int3_appC} as follows
\beq\label{eq:first_int4_appC}
\begin{alignedat}{1}
\int_{\left\{\by \in \Rn \mid\, \normtwo{\by} \leqslant r\right\}} (-G(\by)) \diff{\by} &= \int_{\left\{\by \in \Rn \mid\, \normtwo{\by} \leqslant r\right\}} -\left\langle \varphi_J(\by_0|\bx,t),\by-\by_0 \right\rangle e^{-\left(\frac{1}{2t}\left\Vert \bx-\by\right\Vert _{2}^{2} + J(\by)\right)/\epsilon} \diff{\by}  \\
&\leqslant \int_{\left\{\by \in \Rn \mid\, \normtwo{\by} \leqslant r\right\}}  \normtwo{\varphi_J(\by_0|\bx,t)}\normtwo{\by-\by_0}e^{-\left(\frac{1}{2t}\left\Vert \bx-\by\right\Vert _{2}^{2} + J(\by)\right)/\epsilon} \diff{\by}\\
&\leqslant \normtwo{\varphi_J(\by_0|\bx,t)} \int_{\Rn}\normtwo{\by-\by_0}e^{-\left(\frac{1}{2t}\left\Vert \bx-\by\right\Vert _{2}^{2}\right)/\epsilon} \diff{\by} \\
&= \normtwo{\varphi_J(\by_0|\bx,t)}C_1(\bx,\by_0,t,\epsilon),
\end{alignedat}
\eeq
where $C_1(\bx,\by_0,t,\epsilon)$ was defined in~\eqref{eq:c1_appC}. Combine~\eqref{eq:first_int3_appC} and~\eqref{eq:first_int4_appC} to find
\beq\label{eq:first_int5_appC}
\begin{alignedat}{1}
0 \leqslant \int_{\Rn} F(\by)-G(\by) \diff{\by} &\leqslant \lim_{r \to +\infty} \left(\int_{\left\{\by \in \Rn \mid\, \normtwo{\by} \leqslant r\right\}} F(\by) \diff{\by} + \normtwo{\varphi_J(\by_0|\bx,t)}C_1(\bx,\by_0,t,\epsilon)\right) \\
&= \left(\lim_{r \to +\infty} \int_{\left\{\by \in \Rn \mid\, \normtwo{\by} \leqslant r\right\}} F(\by) \diff{\by}\right)  + \normtwo{\varphi_J(\by_0|\bx,t)}C_1(\bx,\by_0,t,\epsilon).
\end{alignedat}
\eeq
The integral on the right hand side of~\eqref{eq:first_int5_appC} can be bounded using assumption (A3) as follows
\beq\label{eq:first_int6_appC}
\begin{alignedat}{1}
\int_{\left\{\by \in \Rn \mid\, \normtwo{\by} \leqslant r\right\}} F(\by) \diff{\by} &= \int_{\left\{\by \in \Rn \mid\, \normtwo{\by} \leqslant r\right\}} \left\langle \varphi_J(\by|\bx,t),\by-\by_0 \right\rangle e^{-\left(\frac{1}{2t}\left\Vert \bx-\by\right\Vert _{2}^{2} + J(\by)\right)/\epsilon} \diff{\by} \\
&= \int_{\left\{\by \in \Rn \mid\, \normtwo{\by} \leqslant r\right\}} (\left\langle \varphi_J(\by|\bx,t),\by-\by_0 \right\rangle +(n\epsilon - n\epsilon)) e^{-\left(\frac{1}{2t}\left\Vert \bx-\by\right\Vert _{2}^{2} + J(\by)\right)/\epsilon} \diff{\by} \\
&= \int_{\left\{\by \in \Rn \mid\, \normtwo{\by} \leqslant r\right\}} (\left\langle \varphi_J(\by|\bx,t),\by-\by_0 \right\rangle - n\epsilon) e^{-\left(\frac{1}{2t}\left\Vert \bx-\by\right\Vert _{2}^{2} + J(\by)\right)/\epsilon} \diff{\by} \\
&\quad + n\epsilon \int_{\left\{\by \in \Rn \mid\, \normtwo{\by} \leqslant r\right\}} e^{-\left(\frac{1}{2t}\left\Vert \bx-\by\right\Vert _{2}^{2} + J(\by)\right)/\epsilon} \diff{\by} \\
&\leqslant \int_{\left\{\by \in \Rn \mid\, \normtwo{\by} \leqslant r\right\}} (\left\langle \varphi_J(\by|\bx,t),\by-\by_0 \right\rangle - n\epsilon) e^{-\left(\frac{1}{2t}\left\Vert \bx-\by\right\Vert _{2}^{2} + J(\by)\right)/\epsilon} \diff{\by} \\
&\quad + n\epsilon \int_{\Rn} e^{-\left(\frac{1}{2t}\left\Vert \bx-\by\right\Vert _{2}^{2}\right)/\epsilon} \diff{\by} \\
&= \int_{\left\{\by \in \Rn \mid\, \normtwo{\by} \leqslant r\right\}} (\left\langle \varphi_J(\by|\bx,t),\by-\by_0 \right\rangle - n\epsilon) e^{-\left(\frac{1}{2t}\left\Vert \bx-\by\right\Vert _{2}^{2} + J(\by)\right)/\epsilon} \diff{\by} + n\epsilon(2\pi t\epsilon)^{n/2}. \\
\end{alignedat}
\eeq
Combine~\eqref{eq:first_int5_appC} and~\eqref{eq:first_int6_appC} to get
\beq\label{eq:first_int7_appC}
\begin{alignedat}{1}
0 \leqslant \int_{\Rn} F(\by)-G(\by) \diff{\by} &\leqslant \left(\lim_{r \to +\infty}\int_{\left\{\by \in \Rn \mid\, \normtwo{\by} \leqslant r\right\}} (\left\langle \varphi_J(\by|\bx,t),\by-\by_0 \right\rangle - n\epsilon) e^{-\left(\frac{1}{2t}\left\Vert \bx-\by\right\Vert _{2}^{2} + J(\by)\right)/\epsilon} \diff{\by}\right) \\
&\quad+ \normtwo{\varphi_J(\by_0|\bx,t)}C_1(\bx,\by_0,t,\epsilon) + n\epsilon(2\pi t\epsilon)^{n/2}.
\end{alignedat}
\eeq
Combine~\eqref{eq:first_int2_appC} and \eqref{eq:first_int7_appC} to get
\beq\label{eq:first_int8_appC}
\begin{alignedat}{1}
0 \leqslant \int_{\Rn} F(\by)-G(\by) \diff{\by} &= \int_{\Rn} \left\langle \varphi_J(\by|\bx,t)-\varphi_J(\by_0|\bx,t),\by-\by_0 \right\rangle e^{-\left(\frac{1}{2t}\left\Vert \bx-\by\right\Vert _{2}^{2} + J(\by)\right)/\epsilon} \diff{\by} \\
&\leqslant \normtwo{\varphi_J(\by_0|\bx,t)}C_1(\bx,\by_0,t,\epsilon) + n\epsilon(2\pi t\epsilon)^{n/2}.
\end{alignedat}
\eeq
Divide~\eqref{eq:first_int8_appC} by the partition function $Z_J(\bx,t,\epsilon)$ (see Eq.~\eqref{eq:partition_function}) to get 
\beq\label{eq:c2_bound}
0 \leqslant\expectationJ{\left\langle \varphi_J(\by|\bx,t)-\varphi_J(\by_0|\bx,t),\by-\by_0 \right\rangle} \leqslant \frac{\normtwo{\varphi_J(\by_0|\bx,t)}C_1(\bx,\by_0,t,\epsilon) + n\epsilon(2\pi t\epsilon)^{n/2}}{Z_J(\bx,t,\epsilon)} < +\infty.
\eeq
Now, using the Cauchy--Schwarz inequality and~\eqref{eq:norm_appC}, we can bound $\expectationJ{\left|\left\langle\varphi_J(\by_0|\bx,t),\by-\by_0 \right\rangle\right|}$ as follows
\beq\label{eq:c1_bound}
\begin{alignedat}{1}
\expectationJ{\left|\left\langle\varphi_J(\by_0|\bx,t),\by-\by_0 \right\rangle\right|} &= \frac{1}{Z_J(\bx,t,\epsilon)}\int_{\Rn} \left|\left\langle \varphi_J(\by_0|\bx,t),\by-\by_0 \right\rangle\right| e^{-\left(\frac{1}{2t}\left\Vert \bx-\by\right\Vert _{2}^{2} + J(\by)\right)/\epsilon} \diff{\by}\\
&\leqslant \normtwo{\varphi_J(\by_0|\bx,t)}\frac{1}{Z_J(\bx,t,\epsilon)} \int_{\Rn}\normtwo{\by-\by_0}e^{-\left(\frac{1}{2t}\left\Vert \bx-\by\right\Vert _{2}^{2}\right)/\epsilon} \diff{\by} \\
&= \frac{\normtwo{\varphi_J(\by_0|\bx,t)}C_1(\bx,\by_0,t,\epsilon)}{Z_J(\bx,t,\epsilon)}.
\end{alignedat}
\eeq
Use the triangle inequality and the upper bounds in~\eqref{eq:c2_bound} and~\eqref{eq:c1_bound} to obtain
\beq\label{eq:varphi_f_1}
\begin{alignedat}{1}
\expectationJ{\left|\left\langle \varphi_J(\by|\bx,t),\by-\by_0 \right\rangle\right|} &= \expectationJ{\left|\left\langle \varphi_J(\by|\bx,t)-(\varphi_J(\by_0|\bx,t) - \varphi_J(\by_0|\bx,t)),\by-\by_0 \right\rangle\right|} \\
&\leqslant \expectationJ{\left|\left\langle \varphi_J(\by_0|\bx,t),\by-\by_0 \right\rangle\right| + \left|\left\langle \varphi_J(\by|\bx,t)-\varphi_J(\by_0|\bx,t),\by-\by_0 \right\rangle\right|} \\
&= \expectationJ{\left|\left\langle \varphi_J(\by_0|\bx,t),\by-\by_0 \right\rangle\right|} + \expectationJ{\left|\left\langle \varphi_J(\by|\bx,t)-\varphi_J(\by_0|\bx,t),\by-\by_0 \right\rangle\right|} \\
&\leqslant \frac{\normtwo{\varphi_J(\by_0|\bx,t)}C_1(\bx,\by_0,t,\epsilon)}{Z_J(\bx,t,\epsilon)} + \frac{\normtwo{\varphi_J(\by_0|\bx,t)}C_1(\bx,\by_0,t,\epsilon) + n\epsilon(2\pi t\epsilon)^{n/2}}{Z_J(\bx,t,\epsilon)} \\
&<+\infty.
\end{alignedat}
\eeq
Since $\expectationJ{\left|\left\langle \varphi_J(\by|\bx,t),\by-\by_0 \right\rangle\right|} < +\infty$, we can use~\eqref{eq:first_int2_appC} to conclude that
\beq\label{eq:first_int9_appC}
\begin{alignedat}{1}
\expectationJ{\left\langle \varphi_J(\by|\bx,t),\by-\by_0 \right\rangle} &= \frac{1}{Z_J(\bx,t,\epsilon)}\int_{\Rn} \left\langle \varphi_J(\by|\bx,t),\by-\by_0 \right\rangle e^{-\left(\frac{1}{2t}\left\Vert \bx-\by\right\Vert _{2}^{2}+J(\by)\right)/\epsilon} \diff{\by} \\
&= \frac{1}{Z_J(\bx,t,\epsilon)}\lim_{r \to +\infty}\int_{\left\{\by \in \Rn \mid\, \normtwo{\by} \leqslant r\right\}} \left\langle \varphi_J(\by|\bx,t),\by-\by_0 \right\rangle e^{-\left(\frac{1}{2t}\left\Vert \bx-\by\right\Vert _{2}^{2}+J(\by)\right)/\epsilon} \diff{\by} \\
&= \frac{1}{Z_J(\bx,t,\epsilon)}\lim_{r \to +\infty}\int_{\left\{\by \in \Rn \mid\, \normtwo{\by} \leqslant r\right\}} (n\epsilon - n\epsilon + \left\langle \varphi_J(\by|\bx,t),\by-\by_0 \right\rangle) e^{-\left(\frac{1}{2t}\left\Vert \bx-\by\right\Vert _{2}^{2}+J(\by)\right)/\epsilon} \diff{\by} \\
&= n\epsilon - \lim_{r \to +\infty}\left(\int_{\left\{\by \in \Rn \mid\, \normtwo{\by} \leqslant r\right\}} (n\epsilon - \left\langle \varphi_J(\by|\bx,t),\by-\by_0 \right\rangle) e^{-\left(\frac{1}{2t}\left\Vert \bx-\by\right\Vert _{2}^{2}+J(\by)\right)/\epsilon} \diff{\by}\right) \\
&= n\epsilon.
\end{alignedat}
\eeq
Inequality~\eqref{eq:varphi_f_1} and equality~\eqref{eq:first_int9_appC} show the desired results $\expectationJ{\left|\left\langle \varphi_J(\by|\bx,t),\by-\by_0 \right\rangle\right|} < +\infty$ and $\expectationJ{\left\langle \varphi_J(\by|\bx,t),\by-\by_0 \right\rangle} = n\epsilon$, which, after recalling the definition $\varphi_J(\by|\bx,t) = \left(\frac{\by-\bx}{t}\right) + \pi_{\partial J(\by)}(\boldsymbol{0})$, also proves formula~\eqref{eq:diff_mono_rep}.

%%%%%%%%%%%%%%%%%%%%%%%%%%%%%%%%



%%%%%%%%%%%%%%%%%%%%%%%%%%%%%%%%
\bigbreak
\noindent
\textbf{Step 3.} Thanks to Step 2, we have $\expectationJ{\left|\left\langle \varphi_J(\by|\bx,t),\by-\by_0 \right\rangle\right|} < +\infty$ and $\expectationJ{\left\langle \varphi_J(\by|\bx,t),\by-\by_0 \right\rangle} = n\epsilon$ for every $\by_0 \in \Rn$. In particular, the choice of $\by_0 = \boldsymbol{0}$ yields $\expectationJ{\left|\left\langle \varphi_J(\by|\bx,t),\by \right\rangle\right|} < +\infty$ and $\expectationJ{\left\langle \varphi_J(\by|\bx,t),\by \right\rangle} = n\epsilon$. As a consequence, we have that
\beq\label{eq:yzero_1}
\begin{alignedat}{1}
\expectationJ{\left|\left\langle \varphi_J(\by|\bx,t),\by_0 \right\rangle\right|} &= \expectationJ{\left|\left\langle \varphi_J(\by|\bx,t),\by_0 + (\by-\by)\right\rangle\right|} \\
&\leqslant \expectationJ{\left|\left\langle \varphi_J(\by|\bx,t),\by -\by_0 \right\rangle\right| + \left|\left\langle \varphi_J(\by|\bx,t),\by \right\rangle\right|} \\
&= \expectationJ{\left|\left\langle \varphi_J(\by|\bx,t),\by -\by_0 \right\rangle\right|} + \expectationJ{\left|\left\langle \varphi_J(\by|\bx,t),\by \right\rangle\right|} \\
&<\infty,
\end{alignedat}
\eeq
and 
\beq\label{eq:varphi_zero}
\begin{alignedat}{1}
\expectationJ{\left\langle \varphi_J(\by|\bx,t),\by_0 \right\rangle} &= \expectationJ{\left\langle \varphi_J(\by|\bx,t),(\by - \by) + \by_0 \right\rangle}\\
&= \expectationJ{\left\langle \varphi_J(\by|\bx,t),\by \right\rangle} - \expectationJ{\left\langle \varphi_J(\by|\bx,t),\by - \by_0 \right\rangle}\\
&= n\epsilon - n\epsilon \\
&= 0,
\end{alignedat}
\eeq
for every $\by_0 \in \Rn$. Now let $\{\be_i\}_{i=1}^{n}$ denote the standard basis in $\Rn$ and let $\{\varphi_J(\by|\bx,t)_i\}_{i=1}^{n}$ denote the components of the vector $\varphi_J(\by|\bx,t)$, i.e., $\varphi_J(\by|\bx,t) = (\varphi_J(\by|\bx,t)_1,\dots,\varphi_J(\by|\bx,t)_n)$. Using~\eqref{eq:yzero_1} with the choice of $\by_0 = \be_i$ for $i\in\{1,\dots,n\}$, we get $\expectationJ{|\varphi_J(\by|\bx,t)_i|} < +\infty$ for every $i \in \{1,\dots,n\}$. Using the norm inequality $\normtwo{\varphi_J(\by|\bx,t)} \leqslant \sum_{i=1}^{n}\left|\varphi_J(\by|\bx,t)_i\right|$, we can bound $\expectationJ{\normtwo{\varphi_J(\by|\bx,t)}}$ as follows
\beq\label{eq:varphi_bound}
\begin{alignedat}{1}
\expectationJ{\normtwo{\varphi_J(\by|\bx,t)}} &\leqslant \expectationJ{\sum_{i=1}^{n}\left|\varphi_J(\by|\bx,t)_i\right|} \\
&= \sum_{i=1}^{n} \expectationJ{\left|\varphi_J(\by|\bx,t)_i\right|} \\
&< +\infty.
\end{alignedat}
\eeq
We can therefore combine~\eqref{eq:varphi_zero} and~\eqref{eq:varphi_bound} to get $\expectationJ{\left\langle \varphi_J(\by|\bx,t),\by_0 \right\rangle} = \left\langle \expectationJ{\varphi_J(\by|\bx,t)},\by_0 \right\rangle = 0$ for every $\by_0 \in \Rn$, which yields the following equality:
\beq\label{eq:varphi_alone_zero}
\expectationJ{\varphi_J(\by|\bx,t)} = 0.
\eeq
Moreover, recalling the definition $\varphi_J(\by|\bx,t) = \left(\frac{\by-\bx}{t}\right) + \pi_{\partial J(\by)}(\boldsymbol{0})$ and using~\eqref{eq:norm_appC} (with $\by_0 = \bx$) and~\eqref{eq:varphi_bound}, we can bound $\expectationJ{\normtwo{\pi_{\partial J(\by)}(\boldsymbol{0})}}$ as follows
\beq\label{eq:mss_bound_appC}
\begin{alignedat}{1}
\expectationJ{\normtwo{\pi_{\partial J(\by)}(\boldsymbol{0})}} &= \expectationJ{\normtwo{\pi_{\partial J(\by)}(\boldsymbol{0}) + \left(\frac{\by-\bx}{t}\right) - \left(\frac{\by-\bx}{t}\right)}} \\
&\leqslant \expectationJ{\normtwo{\pi_{\partial J(\by)}(\boldsymbol{0}) + \left(\frac{\by-\bx}{t}\right)} + \normtwo{\left(\frac{\by-\bx}{t}\right)}} \\
&=\expectationJ{\normtwo{\varphi_J(\by|\bx,t)}} + \frac{1}{t}\expectationJ{\normtwo{\by-\bx}} \\
&\leqslant \expectationJ{\normtwo{\varphi_J(\by|\bx,t)}} + \frac{C_1(\bx,\bx,t,\epsilon)}{tZ_J(\bx,t,\epsilon)} \\
& < +\infty.
\end{alignedat}
\eeq
We can now combine~\eqref{eq:norm_appC},~\eqref{eq:varphi_alone_zero} and~\eqref{eq:mss_bound_appC} to expand the expected value of $\expectationJ{\varphi_J(\by|\bx,t)}$ as follows
\beq\label{eq:first_int10_appC}
\begin{alignedat}{1}
\expectationJ{\varphi_J(\by|\bx,t)} = \expectationJ{\left(\frac{\by-\bx}{t}\right) + \pi_{\partial J(\by)}(\boldsymbol{0})} &= \expectationJ{\left(\frac{\by-\bx}{t}\right)} + \expectationJ{\pi_{\partial J(\by)}(\boldsymbol{0})} \\
&= \left(\frac{\bu_{PM}(\bx,t,\epsilon)-\bx}{t}\right) + \expectationJ{\pi_{\partial J(\by)}(\boldsymbol{0})} \\
&=\boldsymbol{0}.
\end{alignedat}
\eeq
Solving for $\bu_{PM}(\bx,t,\epsilon)$ in~\eqref{eq:first_int10_appC} yields $\bu_{PM}(\bx,t,\epsilon) = \bx -t\expectationJ{\pi_{\partial J(\by)}(\boldsymbol{0})}$, which gives the representation formula~\eqref{eq:diff_pm_rep}.

We now derive the second representation formula~\eqref{eq:diff_var_rep}. Let $\by_0 = \bu_{PM}(\bx,t,\epsilon)$ in Eq.~\eqref{eq:first_int9_appC} and use the representation formula~\eqref{eq:diff_pm_rep} to find
\beq \label{eq:first_int11_appC}
\begin{alignedat}{1}
\expectationJ{\left\langle \pi_{\partial J(\by)}(\boldsymbol{0}), \by - \bu_{PM}(\bx,t,\epsilon) \right\rangle} &= \expectationJ{\left\langle \pi_{\partial J(\by)}(\boldsymbol{0}) + \left(\frac{\by-\bx}{t}\right) - \left(\frac{\by-\bx}{t}\right), \by - \bu_{PM}(\bx,t,\epsilon) \right\rangle}\\
&= \expectationJ{\left\langle \pi_{\partial J(\by)}(\boldsymbol{0}) + \left(\frac{\by-\bx}{t}\right), \by - \bu_{PM}(\bx,t,\epsilon) \right\rangle - \left\langle \left(\frac{\by-\bx}{t}\right), \by - \bu_{PM}(\bx,t,\epsilon) \right\rangle} \\
&= \expectationJ{\left\langle \pi_{\partial J(\by)}(\boldsymbol{0}) + \left(\frac{\by-\bx}{t}\right), \by - \bu_{PM}(\bx,t,\epsilon) \right\rangle} - \expectationJ{\left\langle \left(\frac{\by-\bx}{t}\right), \by - \bu_{PM}(\bx,t,\epsilon) \right\rangle} \\
&= \expectationJ{\left\langle \varphi_J(\by|\bx,t), \by - \bu_{PM}(\bx,t,\epsilon) \right\rangle} - \expectationJ{\left\langle \left(\frac{\by-\bx}{t}\right), \by - \bu_{PM}(\bx,t,\epsilon) \right\rangle}\\
&= n\epsilon - \expectationJ{\left\langle \left(\frac{\by-\bx}{t}\right), \by - \bu_{PM}(\bx,t,\epsilon) \right\rangle}.
\end{alignedat}
\eeq
We will use~\eqref{eq:first_int11_appC} to derive a representation formula for $\expectationJ{\normsq{\by - \bu_{PM}(\bx,t,\epsilon)}}$. Multiply~\eqref{eq:first_int11_appC} by $t$ and rearrange to get
\beq\label{eq:first_int12_appC}
\expectationJ{\left\langle \by - \bx, \by - \bu_{PM}(\bx,t,\epsilon) \right\rangle} = nt\epsilon - t\expectationJ{\left\langle \pi_{\partial J(\by)}(\boldsymbol{0}), \by - \bu_{PM}(\bx,t,\epsilon) \right\rangle}.
\eeq
The left hand side of~\eqref{eq:first_int12_appC} can be expressed as
\beq\label{eq:rep2_appC}
\begin{alignedat}{1}
\expectationJ{\left\langle \by - \bx, \by - \bu_{PM}(\bx,t,\epsilon) \right\rangle}
&= \expectationJ{\left\langle \by - \bx + (\bu_{PM}(\bx,t,\epsilon) - \bu_{PM}(\bx,t,\epsilon)), \by - \bu_{PM}(\bx,t,\epsilon) \right\rangle} \\
&= \expectationJ{\left\langle \by - \bu_{PM}(\bx,t,\epsilon), \by - \bu_{PM}(\bx,t,\epsilon) \right\rangle + \left\langle \bu_{PM}(\bx,t,\epsilon), \by - \bu_{PM}(\bx,t,\epsilon) \right\rangle } \\
&=\expectationJ{\normsq{\by-\bu_{PM}(\bx,t,\epsilon)}} + \expectationJ{\left\langle\bu_{PM}(\bx,t,\epsilon),\by-\bu_{PM}(\bx,t,\epsilon)\right\rangle} \\
&= \expectationJ{\normsq{\by-\bu_{PM}(\bx,t,\epsilon)}} + \left\langle\bu_{PM}(\bx,t,\epsilon),\expectationJ{\by}-\bu_{PM}(\bx,t,\epsilon)\right\rangle \\
&= \expectationJ{\normsq{\by-\bu_{PM}(\bx,t,\epsilon)}} + \left\langle\bu_{PM}(\bx,t,\epsilon),\bu_{PM}(\bx,t,\epsilon)-\bu_{PM}(\bx,t,\epsilon)\right\rangle \\
&= \expectationJ{\normsq{\by-\bu_{PM}(\bx,t,\epsilon)}}.
\end{alignedat}
\eeq
Combine Eqs.~\eqref{eq:first_int12_appC} and~\eqref{eq:rep2_appC} to get
\[
\expectationJ{\normsq{\by-\bu_{PM}(\bx,t,\epsilon)}} = nt\epsilon - t\expectationJ{\left\langle \pi_{\partial J(\by)}(\boldsymbol{0}), \by - \bu_{PM}(\bx,t,\epsilon)\right\rangle},
\]
which gives the representation formula~\eqref{eq:diff_var_rep}.


%%%%%%%%%%%%%%%%%%%%%%%%%%%%%%%%



%%%%%%%%%%%%%%%%%%%%%%%%%%%%%%%%
\bigbreak
\noindent
\textbf{Step 4.} Thanks to Step 3, the representation formulas~\eqref{eq:diff_pm_rep} and~\eqref{eq:diff_var_rep} hold. Recall that by Prop.~\ref{thm:visc_hj_eqns}(iii), the gradient $\nabla_{\bx}S_\epsilon(\bx,t)$ and Laplacian $\Laplacian S_\epsilon(\bx,t)$ of the solution $S_\epsilon$ to the viscous HJ PDE~\eqref{eq:hj_visc_pde} satisfy the representation formulas
\begin{equation}\label{eq:step4_1}
    \bu_{PM}(\bx,t,\epsilon) = \bx - t\nabla_{\bx}S_\epsilon(\bx,t)
\end{equation}
and
\begin{equation}\label{eq:step4_2}
    \expectationJ{\left \Vert \by - \bu_{PM}(\bx,t,\epsilon) \right \Vert_{2}^{2}} = nt\epsilon - t^2\epsilon\Laplacian S_{\epsilon}(\bx,t).
\end{equation}
Use~\eqref{eq:diff_pm_rep} and~\eqref{eq:step4_1} to get
\[
\nabla_{\bx}S_\epsilon(\bx,t) = \expectationJ{\pi_{\partial J(\by)}(\boldsymbol{0})},
\]
which is the representation formula~\eqref{eq:diff_seps_rep}. Use~\eqref{eq:diff_var_rep} and~\eqref{eq:step4_2} to get
\[
t^2\epsilon\Laplacian S_{\epsilon}(\bx,t) = t\expectationJ{\left\langle \pi_{\partial J(\by)}(\boldsymbol{0}), \by - \bu_{PM}(\bx,t,\epsilon)\right\rangle}
\]
which is, after dividing by $t\epsilon$ on both sides, the representation formulas~\eqref{eq:diff_laplacian_rep}. This concludes Step 4.

%%%%%%%%%%%%%%%%%%%%%%%%%%%%%%%%%%%%%%%%%%%%%%
\textbf{Proof of (ii):} Here we only assume that $J$ satisfies assumptions (A1)-(A3); we do not assume that $\dom J = \Rn$. Let $\{\mu_k\}_{k=1}^{+\infty}$ be a sequence of positive real numbers converging to zero. Define $f_k\colon \Rn\times(0,+\infty)\times(0,+\infty) \to \R$ by
\beq\label{eq:appC_ii_1}
f_\epsilon(\bx,t,k) = -\epsilon\log\left(\frac{1}{(2\pi t\epsilon)^{n/2}}\int_{\Rn}e^{-\left(\frac{1}{2t}\left\Vert \bx-\by\right\Vert _{2}^{2}+S_0(\by,\mu_k)\right)/\epsilon}\diff\by\right)
\eeq
and let $S_0(\bx,\mu_k)$ denote the solution to the first-order HJ PDE~\eqref{eq:hj_non_visc_pde} with initial data $J$ evaluated at $(\bx,\mu_k)$, that is,
\beq\label{eq:appC_ii_2}
S_0(\bx,\mu_k) = \inf_{\by\in\Rn}\left\{\frac{1}{2\mu_k}\normsq{\bx-\by} + J(\by)\right\}.
\eeq
By Prop.~\ref{thm:e_u_nonvisc_hj}(i), the function $\Rn \ni \bx \mapsto S_0(\bx,\mu_k)$ is continuously differentiable and convex for each $k \in \mathbb{N}$, and the sequence of real numbers $\{S_0(\bx,\mu_k)\}_{k=1}^{+\infty}$ converges to $J(\bx)$ for every $\bx \in \dom J$. Moreover, by assumption (A3) ($\inf_{\by \in \Rn} J(\by) = 0$) the sequence $\{S_0(\bx,\mu_k)\}_{k=1}^{+\infty}$ is uniformly bounded from below by $0$, that is,
\begin{align*}
S_0(\bx,\mu_k) & =\inf_{\by\in\Rn}\left\{\frac{1}{2\mu_k}\normsq{\bx-\by} + J(\by)\right\} \\
& \geqslant\inf_{\by\in\Rn}\left\{\frac{1}{2\mu_k}\normsq{\bx-\by}\right\} +\inf_{\by\in\Rn}J(\by)\\
& = 0.
\end{align*}
As a consequence, we can invoke Prop.~\ref{thm:visc_hj_eqns}(i) to conclude that for each $k \in \mathbb{N}$, the function $(\bx,t) \mapsto f_\epsilon(\bx,t,k)$ corresponds to the solution to the viscous HJ PDE~\eqref{eq:hj_visc_pde} with initial data $f_k(\bx,0,\epsilon) = S_0(\bx,\mu_k)$. Moreover, $\Rn \ni \bx \mapsto f_\epsilon(\bx,t,k)$ is continuously differentiable and convex by Prop.~\ref{thm:visc_hj_eqns}(i) and (ii)(a). Finally, as the domain of the function $\bx \mapsto S_0(\bx,\mu_k)$ is $\Rn$, we can use the representation formula~\eqref{eq:diff_seps_rep} in Prop.~\ref{prop:pm_rep_props}(i) (which was proven previously in this Appendix) to express the gradient $\nabla_{\bx}f_{k}(\bx,t,\epsilon)$ as follows
\beq\label{eq:appC_ii_3}
\nabla_{\bx}f_\epsilon(\bx,t,k) = \frac{\int_{\Rn} \nabla_{\by}S_0(\by,\mu_k)e^{-\left(\frac{1}{2t}\left\Vert \bx-\by\right\Vert _{2}^{2}+S_0(\by,\mu_k)\right)/\epsilon} \diff\by}{\int_{\Rn} e^{-\left(\frac{1}{2t}\left\Vert \bx-\by\right\Vert _{2}^{2}+S_0(\by,\mu_k)\right)/\epsilon} \diff\by}.
\eeq

Now, since $S_0(\bx,\mu_k) \geqslant 0$ for every $k \in \mathbb{N}$, we can bound the integrand in~\eqref{eq:appC_ii_1} as follows
\begin{equation}\label{eq:bound_dct_appC}
\frac{1}{(2\pi t\epsilon)^{n/2}}e^{-\left(\frac{1}{2t}\left\Vert \bx-\by\right\Vert _{2}^{2}+S_0(\by,\mu_k)\right)/\epsilon} \leqslant\frac{1}{(2\pi t\epsilon)^{n/2}}e^{-\frac{1}{2t\epsilon}\left\Vert \bx-\by\right\Vert _{2}^{2}},
\end{equation}
where $\int_{\Rn} \frac{1}{(2\pi t\epsilon)^{n/2}}e^{-\frac{1}{2t\epsilon}\left\Vert \bx-\by\right\Vert _{2}^{2}} \diff{\by} = 1$. We can therefore invoke the Lebesgue dominated convergence theorem (\citep{folland2013real}, Theorem 2.24) and use~\eqref{eq:appC_ii_1} and the limit $\lim_{k\to+\infty}e^{-S_0(\bx,\mu_k)/\epsilon}=e^{-J(\bx)/\epsilon}$ (with $\lim_{k\to+\infty}e^{-S_0(\bx,\mu_k)/\epsilon}=0$ for every $\bx\notin\dom J$) to find
\beq\label{eq:lim_final_appC}
\begin{alignedat}{1}
\lim_{k\to+\infty} f_\epsilon(\bx,t,k) &= \lim_{k\to+\infty} -\epsilon\log\left(\frac{1}{(2\pi t\epsilon)^{n/2}}\int_{\Rn}e^{-\left(\frac{1}{2t}\left\Vert \bx-\by\right\Vert _{2}^{2}+S_0(\by,\mu_k)\right)/\epsilon}\diff\by\right) \\
&= -\epsilon\log\left(\frac{1}{(2\pi t\epsilon)^{n/2}}\int_{\dom J}e^{-\left(\frac{1}{2t}\left\Vert \bx-\by\right\Vert _{2}^{2}+J(\by)\right)/\epsilon}\diff\by\right) \\
&= S_\epsilon(\bx,t),
\end{alignedat}
\eeq
which gives the limit~\eqref{eq:sol_approx_limit}. By continuous differentiability and convexity of $\Rn \ni \bx\mapsto f_\epsilon(\bx,t,k)$ and $\Rn \ni \bx \mapsto S_\epsilon(\bx,t)$ and the~limit~\eqref{eq:lim_final_appC}, we can invoke (\cite{rockafellar1970convex}, Theorem 25.7) to conclude that the gradient $\nabla_{\bx}f_k(\bx,t,\mu_k)$ converges to the gradient $\nabla_{\bx}S_\epsilon(\bx,t)$ as $k \to +\infty$. Hence we can take the limit $k\to +\infty$ in~\eqref{eq:appC_ii_3} to find
\beq\label{eq:appC_ii_4}
\begin{alignedat}{1}
\lim_{k\to +\infty} \nabla_{\bx}f_\epsilon(\bx,t,k) &= \lim_{k\to +\infty} \left(\frac{\int_{\Rn} \nabla_{\by}S_0(\by,\mu_k)e^{-\left(\frac{1}{2t}\left\Vert \bx-\by\right\Vert _{2}^{2}+S_0(\by,\mu_k)\right)/\epsilon} \diff\by}{\int_{\Rn} e^{-\left(\frac{1}{2t}\left\Vert \bx-\by\right\Vert _{2}^{2}+S_0(\by,\mu_k)\right)/\epsilon} \diff\by}\right) \\
&= \nabla_{\bx}S_\epsilon(\bx,t),
\end{alignedat}
\eeq
which gives the limit~\eqref{eq:seps_approx_limit}. Finally, using the definition of the posterior mean estimate~\eqref{eq:pm_estimate},~\eqref{eq:appC_ii_4}, and the representation formula~\eqref{eq:grad_s_eps} derived in Prop.~\ref{thm:visc_hj_eqns}(iii), namely $\bu_{PM}(\bx,t,\epsilon) = \bx - t\nabla_{\bx}S_\epsilon(\bx,t)$, we find the two limits
\[
\begin{alignedat}{1}
\bu_{PM}(\bx,t,\epsilon) &= \lim_{k\to +\infty} \left(\frac{\int_{\Rn} \by e^{-\left(\frac{1}{2t}\left\Vert \bx-\by\right\Vert _{2}^{2}+S_0(\by,\mu_k)\right)/\epsilon} \diff\by}{\int_{\Rn} e^{-\left(\frac{1}{2t}\left\Vert \bx-\by\right\Vert _{2}^{2}+S_0(\by,\mu_k)\right)/\epsilon} \diff\by}\right)\\
&= \bx - t\lim_{k\to +\infty} \left(\frac{\int_{\Rn} \nabla_{\by}S_0(\by,\mu_k)e^{-\left(\frac{1}{2t}\left\Vert \bx-\by\right\Vert _{2}^{2}+S_0(\by,\mu_k)\right)/\epsilon} \diff\by}{\int_{\Rn} e^{-\left(\frac{1}{2t}\left\Vert \bx-\by\right\Vert _{2}^{2}+S_0(\by,\mu_k)\right)/\epsilon} \diff\by}\right),
\end{alignedat}
\]
which establishes~\eqref{eq:approx_limit}. This concludes the proof of (ii).